\subsection{Reducing winner queries to single-color omissions}
\label{app:core-pure-interpolation}

Suppose we ask for a winner $k$ at a valuation $v$, subject to a list
of forbidden source bits. Both the positive request and a loss with
several posting colors can be expressed using only single-color losses.
The change of inputs matters: different occurrences of the same input
can be blocked by different owners. Throughout, a word is written using
primitive owned choices, expanding any derived union.

First put $H=\{j:k\in v_j\}$. Give $j\in H$ one copy $j^*$;
give $j\notin H$ a copy $(j,s)$ for each owner $s$ such that
$v_j\ne\{s\}$. Let $D_r$ consist of all copies except those tagged
$r$. A relabeling assigns each leaf a copy of its own input independently.

\begin{lemma}[A winner as simultaneous singleton tests]
\label{core:lem:winner-on-copies}
A word $t$ supplies $k$ at $v$ exactly when some relabeling of $t$
wins the singleton tests $r:D_r$ for all owners $r$.
More generally, for $P\subseteq I$ a relabeling winning all tests
with $r\in P$ exists exactly when
\[
 k\in t(v)\quad\text{or}\quad
 \exists s\notin P:\ t(v)\ne\{s\}.
\]
\end{lemma}
\begin{proof}
Write $B_P(X)$ for the displayed condition with value $X$. At a leaf it
is exactly the condition for an available copy to lie in every required
$D_r$. At an $r$-gate one child must pass every required test, the other
every test except possibly $r$. Indeed the $r$-profile takes disjunction
and the other profiles take conjunction. The same rule holds for
$B_P(X\star_r Y)$. If $X=\{r\}$, the output is $X$, and
$B_{P\setminus\{r\}}(X)=B_P(X)$. Every $Y\ne\{r\}$ passes
$B_{P\setminus\{r\}}$ using $s=r$. The case $Y=\{r\}$ is symmetric.
If neither operand equals $\{r\}$, both pass the reduced test and
their union passes $B_P$ precisely when at least one of them does.
Induction proves the claim. Take $P=I$ for the first assertion.
\end{proof}

\begin{theorem}[Reduction to single-color omissions]
\label{core:thm:pure-omission-reduction}
Let $K\ge3$. Given a winner request and finitely many forbidden source
bits, one can construct a list of single-color omissions on copies of
the inputs. A given word meets the original requirements if and only
if some relabeling of its leaves meets the new list. Identifying the
copies recovers that very word.
\end{theorem}
\begin{proof}
Write $Z_j$ for the forbidden colorings with head $j$, and $z_r$ for
the inputs posted at $r$ in $z$. First give $j$ one copy $(j,f)$ for
each function choosing, for every $z\in Z_j$, a color $f(z)$ with
$z_{f(z)}\ne\varnothing$. For each such copy and each $z\in Z_j$,
forbid its source bit when all copies above $z_{f(z)}$ are posted at
$f(z)$ and every other passive input is uncolored. An empty coloring
in $Z_j$ gives $j$ no copies; a word satisfying that loss cannot use
$j$. An empty $Z_j$ gives $j$ one copy.

To see why this works, trace one occurrence of the head from its leaf
to the root. It supplies the mark at a source probe exactly when no
gate on this path has a sibling returning its owner's foreign singleton.
This follows from the definition of owned choice at each gate. Whether
a sibling returns $\{r\}$ depends only on which leaves are posted at
$r$, by singleton semantics. Thus erasing every other posting preserves
an $r$-block. If the word omits a mixed row, every head occurrence has
such a block; choose its color independently at each row to obtain the
tag $f$ of that occurrence. The same blocks enforce all the new losses.
Conversely, if the tagged word omits the prescribed pure row of a
copy, every occurrence of that copy has a block there. Its owner must
be the sole posting color. Restoring the other colors and identifying
the copies preserves that owner's singleton test and hence the block.

Pull back $v$ to these copies and apply \cref{core:lem:winner-on-copies}.
Lift each pure omission to every copy of its head, pulling its posted
set back along the label map. This enforces exactly the old omission:
an occurrence of any copy has the same singleton tests along its path,
and every occurrence over the old head must be blocked. Finally each
$r:D_r$ is itself a list of pure omissions, one at each head outside
$D_r$ with all of $D_r$ posted at $r$
(\cref{core:eq:total-monochromatic-source}). Both relabelings preserve
the underlying word. The number of inputs in their composite is at most
\[
 K\sum_j\prod_{z\in Z_j}|\{r:z_r\ne\varnothing\}|.
\]
\end{proof}

For a full guard table one can apply the reduction to each maximal
positive bit with the same minimal zeros. Taking the union of the
resulting witnesses after identifying their inputs preserves all the
zeros and supplies every positive, by antitonicity and source
decomposition. As usual an identically zero table has no realizer.

\paragraph{The remaining simultaneous constraint.}
For each owner $r$, let $\mathcal F_r$ consist of the subsets $S$ of
the new input set $W$ satisfying
\[
  j\in S\ \Longrightarrow\ S\cap T\ne\varnothing
\]
for every pure $r$-omission with head $j$ and posted set $T$. This
family contains $\varnothing$ and is closed under unions. Conversely
every union-closed family $\mathcal F$ containing $\varnothing$ on a
finite $W$ arises this way. For $S\notin\mathcal F$, choose
$j\in S$ outside the union of the feasible subsets of $S$. The
implication with tail $W\setminus S$ excludes $S$ and holds throughout
$\mathcal F$. An empty tail simply forbids the head.

An \emph{$r$-pruning} keeps both children at every $r$-gate it visits
and chooses one child at every other gate. A word meets the pure list
if and only if every $r$-pruning has its set of leaf labels in
$\mathcal F_r$, for every $r$. Indeed at a pure $r$-probe every foreign
gate acts as union. Since $\star_r$ distributes over union, the value is
the union of the $r$-folds on the pruning supports. A fold transmits
the head $j$ precisely when its support contains $j$ and misses $T$.
This gives an equivalent problem on arbitrary union-closed families,
one per owner, with no distinguished winner or mixed row.

\begin{corollary}[One arbitrary family of pruning supports]
\label{core:cor:one-arbitrary-pruning-family}
Fix an owner $a$ and allow its union-closed family $\mathcal F_a$ to
be arbitrary. Suppose that at every other owner the nonempty feasible
supports form an upward-closed family. Describe them as the sets
meeting every member of a list $\mathcal D_r$, including $W$ as a
redundant test. A word is feasible exactly when every selection
$A_r\in\mathcal D_r$, one for each $r\ne a$, has intersection
containing a nonempty member of $\mathcal F_a$.
\end{corollary}
\begin{proof}
Every $r$-pruning meets $A_r$ exactly when the word passes the
$r$-singleton test on $A_r$, by the disjunction and conjunction rules
for profiles. Fix one selection. At each visited $r$-gate with
$r\ne a$, retain a child passing its own test. Every other required
test passes at both children. Retain both at $a$-gates. The result
is an $a$-pruning, all of whose leaves belong to the intersection.

Conversely, place an $a$-fold of the chosen feasible support at each
complete selection. Combine these folds in a tree that branches over
$\mathcal D_r$ by $\star_r$, taking the owners in any fixed order
and writing each finite branch as a binary fold. Every $a$-pruning
selects one of the original feasible supports. At owner $r$, for any
specified $A_r$, every $r$-pruning retains the branch indexed by
$A_r$ when that owner's turn comes, and hence contains a leaf in
$A_r$. It therefore meets every required set at $r$.
\end{proof}

\paragraph{Two arbitrary families.}
For a word using just two owners, the inclusion-minimal pruning supports
of one owner are the minimal transversals of the other owner's supports.
Indeed a coalition meets every $+$-pruning exactly when the $+$-profile
wins there. Choosing a winning child at each $+$-gate gives a
$-$-pruning contained in the coalition; the converse follows by the
same profile recursion. Thus a necessary condition for two prescribed
families is that they contain such a dual pair. If other owners have
no constraints, their gates may be pruned to either child first:
both constrained owners' surviving prunings were already prunings of
the original word.

\begin{proposition}[Dual pruning supports need not have a common word]
\label{core:prop:dual-pruning-obstruction}
On inputs $a,b,c,d,e$, prescribe the union closures of the following
families for owners $+$ and $-$ respectively, allowing the empty union:
\[
 \mathcal A=\{be,ae,ad,bcd\},\qquad
 \mathcal B=\{ab,ace,de\}.
\]
Here $be$ denotes $\{b,e\}$, and similarly for the other sets.
These two lists are each other's minimal transversals. No word has
all its pruning supports in the prescribed families.
\end{proposition}
\begin{proof}
A transversal of $\mathcal B$ missing $a$ needs $b$ and either $e$
or both $c,d$. A transversal containing $a$ still needs $d$ or $e$.
This gives exactly $\mathcal A$ as minimal transversals; transversal
duality gives the reverse assertion.

Every nonempty union from $\mathcal A$ meets $de$, and every nonempty
union from $\mathcal B$ meets $be$. Thus a word must pass the tests
$+:de$ and $-:be$. Its pruning supports must also obey these four
implications, each of which holds on the displayed generators:
\[
 \begin{array}{c|ll}
 + & c\Rightarrow b & d\Rightarrow a\text{ or }c\\
 - & b\Rightarrow a & e\Rightarrow c\text{ or }d .
 \end{array}
\]
Take a shortest two-owner word meeting the two tests and four omissions.
At a $+$-root, both children retain $-:be$ and both $-$-omissions.
Neither child can pass $+:ac$, since $ac$ misses $be$ and winning
coalitions at different owners intersect. Neither can pass $+:b$:
this would give a $-$-pruning inside $\{b\}$, contrary to $b\Rightarrow a$.
Thus neither child can absorb at either $+$-omission. Both children
must already satisfy them, and one retains $+:de$.

At a $-$-root both children retain $+:de$ and the $+$-omissions.
Neither can pass $-:a$, since $a$ misses $de$. Neither can pass
$-:cd$: a nonempty $+$-pruning inside $cd$ cannot contain $c$ by
the first implication, nor then $d$ by the second. Again both children
satisfy all four omissions and one retains both tests. Either root
contradicts shortest length. A leaf passing both tests must be $e$,
which violates the last implication.
\end{proof}
